\documentclass[10pt,twocolumn]{article}
\usepackage[margin=0.62in]{geometry}
\usepackage{amsmath,amssymb,mathtools,bm}
\usepackage{enumitem}
\usepackage[hidelinks]{hyperref}
\setlength{\parindent}{0pt}
\setlength{\parskip}{3pt}
\setlist[itemize]{leftmargin=*,topsep=2pt,itemsep=1pt}
\setlist[enumerate]{leftmargin=*,topsep=2pt,itemsep=1pt}
\newcommand{\R}{\mathbb R}
\newcommand{\C}{\mathbb C}
\newcommand{\Span}{\operatorname{span}}
\newcommand{\Ree}{\operatorname{Re}}
\newcommand{\clip}{\operatorname{clip}}
\newcommand{\argmax}{\operatorname*{arg\,max}}
\newcommand{\argmin}{\operatorname*{arg\,min}}
\newcommand{\diag}{\operatorname{diag}}
\newcommand{\rank}{\operatorname{rank}}
\newcommand{\eps}{\varepsilon}
\newcommand{\lcb}{\mathrm{lcb}}
\newcommand{\plat}{\mathrm{plat}}
\newcommand{\logi}{\mathrm{log}}
\newcommand{\exec}{\mathrm{exec}}
\newcommand{\tr}{\mathrm{TR}}
\newcommand{\vol}{\mathrm{vol}}
\newcommand{\FS}{\mathrm{FS}}
\newcommand{\dd}{\,\mathrm d}
\title{Route C Plateau Acquisition v1.2\\Pedagogical Derivation of the Log-Volume Novelty Rule}
\author{Holstein repository note}
\date{2026-06-08}
\begin{document}
\maketitle
\small
\raggedright

\textbf{Reader contract.} This PDF is the derivation-first companion to
\texttt{MATH/notes/route\_c\_plateau\_acquisition\_v1\_doctrine.md}. The
Markdown file is the repo-agent implementation contract. This PDF explains why
Route C should rank plateau candidates by incremental Fubini--Study log-volume
growth, while preserving the unlock-or-zero-commit energy rule.

\section{What plateau acquisition is trying to measure}

At ADAPT step $t$, the selector owns an admitted ansatz word
\[
\mathcal O_t^{\logi}=(o_1,\ldots,o_{n_t}),\qquad
I_t^{\logi}=\{1,
\ldots,n_t\}.
\]
Some coordinates were bought during plateau mode and currently have committed
amplitude zero. They are not pending candidates. They are admitted coordinates
with provenance
\[
D_t^0=\{j\in I_t^{\logi}:\operatorname{status}_{t,j}=\mathrm{dormant\_plateau}\}.
\]
The executable nonzero circuit is the smaller set
\[
I_t^{\exec}=\{j\in I_t^{\logi}:|\theta_{t,j}|>\theta_{\rm zero}\}.
\]
The physical state is unchanged by zero amplitudes, but the local coordinate
chart has acquired directions. Route C is therefore a chart-growth mechanism:
it buys tangent directions at zero energy cost and later asks the ordinary inner
optimizer whether the enlarged chart contains a validated descent.

Normal Phase III uses
\[
S_3(r;t)=\frac{\Delta E_{\tr}(r;t)\mathcal N_3(r;t)}{1+K_3(r;t)}.
\]
When
\[
\Delta_{\max}^{(3)}(t)=\max_{r\in\mathcal R_3(t)}\Delta E_{\tr}(r;t)
\le \eps_{\plat}(t),
\]
the product score cannot express novelty acquisition, because the gain factor is
unresolved. Plateau mode replaces the selection law; it does not accept an
uphill energy step.

\section{The primitive geometry}

Let $|\psi_t\rangle$ be the committed reportable state. The horizontal tangent
for an admitted coordinate $j$ is
\[
\tau_j(t)=Q_{\psi_t}\partial_{\theta_j}
|\psi(\mathcal O_t^{\logi},\theta_t^{\logi})\rangle .
\]
For a candidate-position record $r=(m,p)$ with trial amplitude $\alpha_r$,
\[
\tau_r(t)=Q_{\psi_t}\partial_{\alpha_r}
|\psi(\mathcal O_t^{\logi}\oplus r,
(\theta_t^{\logi},\alpha_r))\rangle\big|_{\alpha_r=0}.
\]
The plateau context is the admitted scaffold:
\[
C_t=I_t^{\logi}
\]
unless a run explicitly logs a contracted approximation such as
$W_r^{\rm act}\cup D_t^0$. Build
\[
T_C=[\tau_j]_{j\in C_t}.
\]
The real Fubini--Study Gram primitives are
\[
G_C=\Ree(T_C^\dagger T_C),\qquad
q_r=\Ree(T_C^\dagger\tau_r),\qquad
F_r=\Ree\langle\tau_r,\tau_r\rangle .
\]
Componentwise,
\[
(G_C)_{ij}=\Ree\langle\tau_i,\tau_j\rangle,
\quad
(q_r)_i=\Ree\langle\tau_i,\tau_r\rangle,
\quad
F_r=\Ree\langle\tau_r,\tau_r\rangle .
\]
$G_C$ measures metric correlations among already admitted directions, $q_r$
records the candidate's shadows against them, and $F_r$ is the candidate's own
local Fubini--Study length squared.

\section{Projection residual from least squares}

The old span explains $\tau_r$ by synthesizing
\[
T_C\beta=\sum_{j\in C_t}\beta_j\tau_j.
\]
The best old-span explanation solves
\[
\beta^\star\in\argmin_\beta
\|\tau_r-T_C\beta\|_{\FS}^2.
\]
Expand the objective using the real Fubini--Study metric:
\[
\|\tau_r-T_C\beta\|_{\FS}^2
=F_r-2\beta^\top q_r+\beta^\top G_C\beta.
\]
Stationarity gives the normal equation
\[
G_C\beta^\star=q_r,
\qquad
\beta^\star=G_C^+q_r
\]
on the resolved support. Substituting the optimized coefficient eliminates
$\beta$:
\[
\min_\beta\|\tau_r-T_C\beta\|_{\FS}^2
=F_r-q_r^\top G_C^+q_r.
\]
Define the supported residual strength
\[
\sigma_r^\perp=\bigl[F_r-q_r^\top G_C^+q_r\bigr]_+.
\]
The fractional novelty used in earlier doctrine is
\[
\mathcal N_r^{\rm frac}=\clip_{[0,1]}
\frac{\sigma_r^\perp}{\max(F_r,\eps_{\rm red})}.
\]
This answers: what fraction of this candidate's tangent cannot be explained by
already admitted coordinates?

\section{Why the fraction is not enough}

The fraction is scale-free. Scale-free is useful for detecting redundancy, but
plateau acquisition has a second requirement: the added direction must have a
resolved size. Two failure modes follow immediately.

\textbf{Tiny but high fraction.} If $F_r\ll 1$ and the candidate is orthogonal to
$T_C$, then $\mathcal N_r^{\rm frac}\approx1$ while $\sigma_r^\perp\approx F_r$
is numerically useless.

\textbf{Large but mostly redundant.} If $F_r$ is large but almost all of it lies
in $\Span(T_C)$, then the raw candidate is energetic/geometric clutter. Its
fractional residual catches the redundancy.

Thus Route C should gate by both
\[
\sigma_{r,\lcb}^\perp\ge\sigma_{\plat}^{\min},
\qquad
\mathcal N_{r,\lcb}^{\rm frac}\ge\nu_{\plat}^{\min}.
\]
The first gate requires a resolved absolute tangent addition. The second gate
requires that the addition is not mostly old-span motion.

\section{Where the logarithm comes from}

A coordinate chart with tangent matrix $T$ has local Fubini--Study volume element
\[
\dd\operatorname{Vol}_{\FS}(T)
\propto
\sqrt{\det(T^\dagger T)}\,\dd a_1\cdots \dd a_d.
\]
The square root appears because the Gram determinant gives squared volume. If a
new coordinate is appended, the physical infinitesimal volume ratio is therefore
\[
\frac{\dd\operatorname{Vol}_{\FS}(C\cup r)}
{\dd\operatorname{Vol}_{\FS}(C)\,\dd a_r}
=
\sqrt{
\frac{\det G_{C\cup r}}
{\det G_C}}.
\]
This expression assumes nonsingular Grams; the ridge version below is the
numerically supported replacement.
Route C wants an additive acquisition score. Volumes multiply when independent
coordinates are appended, while logarithms turn products into sums:
\[
\log(R_1R_2)=\log R_1+\log R_2.
\]
This is the sole reason for the logarithm: the score is a log-volume increment,
so marginal chart-growth contributions can be compared and added.

For a possibly rank-deficient or noisy Gram, use a ridge
\[
A_C=G_C+\lambda_{\vol}I,
\qquad \lambda_{\vol}>0.
\]
The ridge supplies a baseline metric for the one empty coordinate that would be
present even before measuring the candidate's residual. After appending
candidate $r$, the regularized augmented Gram is
\[
G_{C\cup r}^{(\lambda)}=
\begin{bmatrix}
G_C+\lambda_{\vol}I & q_r\\
q_r^\top & F_r+\lambda_{\vol}
\end{bmatrix}.
\]
The block determinant identity gives
\[
\det G_{C\cup r}^{(\lambda)}
=\det(A_C)
\left(F_r+\lambda_{\vol}-q_r^\top A_C^{-1}q_r\right).
\]
Define the ridge residual
\[
\sigma_{r|C}^{\lambda}=\bigl[F_r-q_r^\top A_C^{-1}q_r\bigr]_+.
\]
Then
\[
F_r+\lambda_{\vol}-q_r^\top A_C^{-1}q_r
=\lambda_{\vol}+\sigma_{r|C}^{\lambda}.
\]
Normalize by the baseline ridge volume of the added coordinate:
\[
R_{r|C}^{\det}
:=
\frac{\det G_{C\cup r}^{(\lambda)}}
{\lambda_{\vol}\det(A_C)}
=1+\frac{\sigma_{r|C}^{\lambda}}{\lambda_{\vol}}.
\]
The physical Fubini--Study volume ratio is the square root:
\[
R_{r|C}^{\FS}=\sqrt{R_{r|C}^{\det}}.
\]
Thus the physical log-volume gain is
\[
\log R_{r|C}^{\FS}
=\frac12\log\left(1+\frac{\sigma_{r|C}^{\lambda}}{\lambda_{\vol}}\right).
\]
The repo score uses the log-determinant version
\[
\Delta V_{r|C}^{\lambda}
:=
\log R_{r|C}^{\det}
=
\log\left(1+\frac{\sigma_{r|C}^{\lambda}}{\lambda_{\vol}}\right).
\]
This is exactly twice the physical log-volume gain. The factor $1/2$ is global
and positive, so it does not change candidate ranking. Thresholds such as
$V_{\plat}^{\min}$ should use the same convention as the implementation.

This is the direct implementation of ``buy the candidate that grows the local
state-manifold chart most.''

\section{Route-C acquisition score}

The v1.2 plateau score is
\[
S_{3,\plat}^{\vol}(r;t)=
\frac{[\Delta V_{r|C}^{\lambda}]_{\lcb}}{1+K_3(r;t)}.
\]
The selector uses
\[
r_t\in\argmax_{r\in\mathcal R_{\plat}(t)}
S_{3,\plat}^{\vol}(r;t),
\]
subject to feasibility, duplicate blocking, and the two residual gates.

This score uses Geo-ADAPT's useful ingredient: the Fubini--Study/QIM metric. It
uses geometry as coverage, not as a natural-gradient energy step. In plateau
mode the immediate energy coordinate is unresolved, so the metric asks which
record expands the admitted tangent chart.

\section{Where Schur response belongs}

A Schur-relaxed tangent often has the form
\[
\tau_r^\star=\tau_r+T_Cs_r.
\]
Let $P_C$ be the Fubini--Study projection onto $\Span(T_C)$. Since $T_Cs_r$ is
already in the context span,
\[
(I-P_C)\tau_r^\star=(I-P_C)(\tau_r+T_Cs_r)=(I-P_C)\tau_r.
\]
Thus if the novelty projector uses the same admitted context, Schur
reparameterization does not change the out-of-context residual. It changes the
in-context coordinate response and the energy model, not the new span.

Route C v1.2 therefore uses raw horizontal tangent geometry for plateau
log-volume acquisition. Schur quantities remain essential for ordinary Phase III
energy ranking, negative-curvature diagnostics, and trial-refit/unlock
validation.

\section{Block-aware plateau acquisition}

Plateau escape can be collective. For a candidate block
$B=\{r_1,\ldots,r_s\}$, define
\[
T_B=[\tau_{r_1},\ldots,\tau_{r_s}],
\quad
F_{BB}=\Ree(T_B^\dagger T_B),
\quad
Q_{CB}=\Ree(T_C^\dagger T_B).
\]
The residual block Gram is
\[
\Sigma_{B|C}^{\lambda}
=F_{BB}-Q_{CB}^\top(G_C+\lambda_{\vol}I)^{-1}Q_{CB}.
\]
The block log-volume increment is
\[
\Delta V_{B|C}^{\lambda}
=\log\det(I_s+\lambda_{\vol}^{-1}\Sigma_{B|C}^{\lambda}).
\]
A practical implementation can remain greedy:
\[
r_{\rm next}\in\argmax_{r\in\mathcal R_{\plat}(t)\setminus B}
\frac{[\Delta V_{r|C\cup B}^{\lambda}]_{\lcb}}{1+K_3(r;t)}.
\]
After one direction has been added to the context, another nearly equivalent
candidate has smaller marginal volume gain. This is the block-aware version of
``do not keep buying the same dormant direction.''

\section{Residual-pool leverage as optional Geo-style coverage}

For plateau candidates $u,v\in\mathcal R_{\plat}(t)$, define
\[
K_{uv|C}^{\lambda}=\Ree\langle\tau_u,\tau_v\rangle
-q_u^\top(G_C+\lambda_{\vol}I)^{-1}q_v.
\]
This is a candidate-pool Gram after subtracting the admitted context. A ridge
leverage score is
\[
\ell_r^{\lambda_P}=
\left[K_{\mathcal R|C}^{\lambda}
(K_{\mathcal R|C}^{\lambda}+\lambda_PI)^{-1}\right]_{rr}.
\]
Leverage can be logged, used as a tie-breaker, or used as a gate. The minimal
v1.2 implementation should not require it. The primary score is log-volume over
cost.

\section{Uncertainty-aware form}

Let
\[
y_r=(G_C+\lambda_{\vol}I)^{-1}q_r.
\]
For
\[
\sigma_{r|C}^{\lambda}=F_r-q_r^\top y_r,
\]
first-order perturbation gives
\[
\dd\sigma_{r|C}^{\lambda}
=\dd F_r-2y_r^\top\dd q_r+y_r^\top(\dd G_C)y_r.
\]
A shot-block bootstrap or covariance propagation can estimate
$\operatorname{SE}(\sigma_{r|C}^{\lambda})$ from the same measured Gram and
overlap primitives. Then
\[
[\Delta V_{r|C}^{\lambda}]_{\lcb}
=\left[\log\left(1+
\frac{[\sigma_{r|C}^{\lambda}-z_\alpha\operatorname{SE}(\sigma_{r|C}^{\lambda})]_+}
{\lambda_{\vol}}
\right)\right]_+.
\]
For deterministic local diagnostics, the LCB can be the raw value. For
measurement-facing routes, low-confidence volume increments should not drive
plateau acquisition.

\section{Trial refit remains ordinary}

After selecting $r_t$, form
\[
\mathcal O_t^{\rm trial}=\mathcal O_t^{\logi}\oplus r_t,
\qquad
x_t^{(0)}=(\theta_t^{\logi},0).
\]
Because $\theta_t^{\logi}$ already contains zeros for dormant coordinates, the
trial vector includes active/base coordinates, dormant coordinates, and the new
candidate coordinate. The ordinary inner optimizer solves
\[
x_t^{\rm opt}\approx\argmin_x E(\mathcal O_t^{\rm trial},x).
\]
For SPSA, every dormant and selected plateau coordinate must appear in the
perturbation vector:
\[
i\in D_t^0\cup\{r_t\}\implies
\Delta_i^{(s)}\in\{-1,+1\},\qquad c_s\Delta_i^{(s)}\ne0.
\]
Finite-amplitude seed probes can help when first-order SPSA is weak at the
origin, but they are optimizer-strength diagnostics. They do not change the
acquisition statistic or the unlock rule.

\section{Energy safety: unlock or zero commit}

Let
\[
E_t^{\rm rep}=\text{current reportable energy},
\qquad
E_t^{\rm trial}=E(\mathcal O_t^{\rm trial},x_t^{\rm opt}).
\]
The raw trial drop is
\[
\widehat\Delta_t^{\rm trial}=E_t^{\rm rep}-E_t^{\rm trial}.
\]
A lower-confidence validated drop is
\[
Y_t^{\rm trial}=\widehat\Delta_t^{\rm trial}
-z_\alpha\sigma_{\Delta E,t}
-b_{E,\Delta}^{\rm cmp}(t)
-\gamma_{\rm hw}\Delta b_{E,r_t}^{\rm inc}(t).
\]
Unlock succeeds if
\[
Y_t^{\rm trial}\ge\eps_{\rm unlock}(t).
\]
On success, commit $x_t^{\rm opt}$ and report $E_t^{\rm trial}$. Plateau-acquired
coordinates whose amplitudes exceed $\theta_{\rm zero}$ become non-dormant.

On failure,
\[
Y_t^{\rm trial}<\eps_{\rm unlock}(t),
\]
commit the selected record logically but restore the energetic state:
\[
E_{t+1}^{\rm rep}=E_t^{\rm rep},
\qquad
\theta_{D_{t+1}^0}=0.
\]
The failed trial's nonzero dormant amplitudes and optimizer memory must not leak
into the next committed state unless a separately named experimental route says
so.

\section{Implementation translation}

The implementation-facing Route C v1.2 selector is:
\[
\Delta_{\max}^{(3)}(t)\le\eps_{\plat}(t)
\implies
r_t\in\argmax_{r\in\mathcal R_{\plat}(t)}
\frac{[\Delta V_{r|I_t^{\logi}}^{\lambda}]_{\lcb}}{1+K_3(r;t)}.
\]
Candidate pass gates are
\[
\sigma_{r,\lcb}^\perp\ge\sigma_{\plat}^{\min},
\qquad
\mathcal N_{r,\lcb}^{\rm frac}\ge\nu_{\plat}^{\min}.
\]
The acquisition context is
\[
I_t^{\logi}=I_t^{\rm base}\cup D_t^0,
\qquad
D_t^0=\{j:\operatorname{status}_{t,j}=\mathrm{dormant\_plateau}\}.
\]
The commit rule remains energy-based:
\begin{align*}
Y_t^{\rm trial}\ge\eps_{\rm unlock}
&\Rightarrow \text{commit nonzero refit},\\
Y_t^{\rm trial}<\eps_{\rm unlock}
&\Rightarrow \text{zero commit}.
\end{align*}

\section{Minimum regression fixtures}

\begin{enumerate}
\item A tangent already in the active-plus-dormant span has
$\sigma^\perp\approx0$ and $\Delta V\approx0$.
\item A tiny orthogonal tangent has high fractional novelty but fails
$\sigma_{\plat}^{\min}$.
\item A large mostly redundant tangent fails $\nu_{\plat}^{\min}$.
\item Among passing candidates, larger $\Delta V/(1+K_3)$ wins.
\item After selecting one member of a redundant pair, the second member's
marginal log-volume gain falls.
\item Dormant coordinates appear in the SPSA vector during trial refit.
\item Failed unlock restores reportable energy and zero-commits plateau
coordinates.
\item A synthetic collective fixture,
\[
E(a_1,a_2)=E_0+\lambda(a_1^2+a_2^2)-\beta a_1a_2,
\quad \beta>2\lambda,
\]
has neutral axes but a descending joint direction.
\end{enumerate}


\section{Worked scalar examples for the gates}

Assume the old span has already been projected out, so the only remaining
quantity is the residual strength. The gates can be read without matrix algebra.

\textbf{Tiny orthogonal direction.} Let $q_r=0$ and $F_r=10^{-12}$. Then
\[
\sigma_r^\perp=10^{-12},
\qquad
\mathcal N_r^{\rm frac}=1.
\]
A fraction-only selector treats this as maximally novel. A plateau chart-growth
selector rejects it if $\sigma_{\plat}^{\min}\gg10^{-12}$. The candidate adds a
new direction in exact arithmetic, but it does not add a resolved state-space
coordinate at the numerical or measurement scale of the run.

\textbf{Large mostly old direction.} Let $F_r=1$ and
$q_r^\top G_C^+q_r=0.9999$. Then
\[
\sigma_r^\perp=10^{-4},
\qquad
\mathcal N_r^{\rm frac}=10^{-4}.
\]
The absolute residual may be visible, but the candidate is almost an old-span
motion. The fractional gate rejects it unless $\nu_{\plat}^{\min}\le10^{-4}$.

\textbf{Resolved and genuinely new.} Let $F_r=10^{-2}$ and the old span explains
only $20\%$ of it. Then
\[
\sigma_r^\perp=8\times10^{-3},
\qquad
\mathcal N_r^{\rm frac}=0.8.
\]
This candidate passes both gates for typical plateau settings and receives a
log-volume rank
\[
\Delta V=\log\left(1+\frac{8\times10^{-3}}{\lambda_{\vol}}\right)
\]
under the ridge approximation.

\section{Why the logarithm is the stable volume unit}

The raw determinant ratio multiplies across independent coordinates. If two
orthogonal residual directions have residual strengths $\sigma_1$ and
$\sigma_2$, the regularized volume ratio is approximately
\[
\left(1+\frac{\sigma_1}{\lambda_{\vol}}\right)
\left(1+\frac{\sigma_2}{\lambda_{\vol}}\right).
\]
Taking logs makes marginal gains additive:
\[
\Delta V_{\{1,2\}}
=\log\left(1+\frac{\sigma_1}{\lambda_{\vol}}\right)
+\log\left(1+\frac{\sigma_2}{\lambda_{\vol}}\right).
\]
This is useful for greedy selection. A candidate already explained by the
current block has small marginal residual, so its added log-volume is small.
The logarithm also compresses extremely large residuals so cost and gates can
still influence ranking.

For small residuals,
\[
\log(1+x)=x+O(x^2),
\qquad x=\sigma/\lambda_{\vol}.
\]
Thus near the resolution floor the log-volume score behaves like a scaled
absolute residual. For large residuals it becomes a diminishing-return chart
coverage score.

\section{Block determinant derivation}

The singleton determinant identity is the $s=1$ case of the block formula. Let
\[
A_C=G_C+\lambda_{\vol}I,
\qquad
B=Q_{CB},
\qquad
D=F_{BB}+\lambda_{\vol}I_s.
\]
The augmented block is
\[
\begin{bmatrix}
A_C & B\\
B^\top & D
\end{bmatrix}.
\]
The Schur determinant identity gives
\[
\det\begin{bmatrix}A_C&B\\B^\top&D\end{bmatrix}
=\det(A_C)\det(D-B^\top A_C^{-1}B).
\]
Substitute $D=F_{BB}+\lambda_{\vol}I_s$:
\[
D-B^\top A_C^{-1}B
=\lambda_{\vol}I_s+
\left(F_{BB}-Q_{CB}^\top A_C^{-1}Q_{CB}\right).
\]
The parenthesized term is
\[
\Sigma_{B|C}^{\lambda}=F_{BB}-Q_{CB}^\top A_C^{-1}Q_{CB}.
\]
The nontrivial volume ratio normalizes out the baseline ridge volume of $s$ new
empty coordinates:
\[
\frac{\det G_{C\cup B}^{(\lambda)}}
{\det(A_C)\lambda_{\vol}^s}
=\det\left(I_s+\lambda_{\vol}^{-1}\Sigma_{B|C}^{\lambda}\right).
\]
Taking logs yields
\[
\Delta V_{B|C}^{\lambda}
=\log\det\left(I_s+\lambda_{\vol}^{-1}\Sigma_{B|C}^{\lambda}\right).
\]
This is the block chart-growth objective. Greedy Route C approximates it by
adding one record at a time with the current block included in the context.

\section{Relation to rank growth}

Rank growth is the hard-threshold version of volume growth. If the eigenvalues
of the residual block are $\mu_1,\ldots,\mu_s$, then
\[
\Delta V_{B|C}^{\lambda}
=\sum_{k=1}^s\log\left(1+\frac{\mu_k}{\lambda_{\vol}}\right).
\]
Resolved numerical rank counts modes above a threshold:
\[
\rank_{\eps_{\rank}}(\Sigma_{B|C})
=\left|\{k:\mu_k>\eps_{\rank}\mu_{\max}\}\right|.
\]
Rank tells how many directions survived the support rule. Log-volume also keeps
track of their sizes. Route C telemetry should ideally report both: rank for
human audit, log-volume for ranking.

\section{How Route C differs from accepting uphill moves}

A plateau admission at zero has exact identity safety:
\[
e^{-i0A_r}=I,
\qquad
U(\mathcal O_t^{\logi}\oplus r, (\theta_t^{\logi},0))
=U(\mathcal O_t^{\logi},\theta_t^{\logi}).
\]
Therefore
\[
E(\mathcal O_t^{\logi}\oplus r, (\theta_t^{\logi},0))
=E_t^{\rm rep}.
\]
The algorithm may evaluate nonzero trial points during refit. Those trial points
are not committed unless the validated drop passes
\[
Y_t^{\rm trial}\ge\eps_{\rm unlock}.
\]
If the trial fails,
\[
E_{t+1}^{\rm rep}=E_t^{\rm rep}.
\]
The new record remains in the logical scaffold at zero amplitude so future
novelty sees it. This is energy-safe chart growth.

\section{Collective escape fixture}

The minimal local-hurdle fixture is
\[
E(a_1,a_2)=E_0+\lambda(a_1^2+a_2^2)-\beta a_1a_2,
\qquad
\lambda>0,
\quad
\beta>2\lambda.
\]
Axis probes do not improve energy:
\[
E(s,0)-E_0=\lambda s^2\ge0,
\qquad
E(0,s)-E_0=\lambda s^2\ge0.
\]
The diagonal direction improves energy:
\[
E(s,s)-E_0=(2\lambda-\beta)s^2<0.
\]
A singleton energy-gain selector can stall because neither coordinate is useful
alone. Route C buys resolved, nonredundant coordinates at zero amplitude until
the ordinary enlarged optimizer has access to the joint direction. The unlock
rule then decides whether the joint move is real.

\section{Finite-angle novelty as fallback}

A finite-angle state probe would compare
\[
|\psi_r(\alpha_0)\rangle
=U(\mathcal O_t^{\logi}\oplus r,(	heta_t^{\logi},\alpha_0))
|\psi_{\rm ref}\rangle
\]
against the current state using
\[
D_{\FS}^2(r;\alpha_0)=1-|\langle\psi_t|\psi_r(\alpha_0)\rangle|^2.
\]
This is a legitimate diagnostic, but it mixes tangent novelty, finite-angle
curvature, ordering effects, and nonlinear state overlap. The default Route C
novelty primitive should remain tangent log-volume. Finite-angle novelty is a
fallback when tangent rank saturates while plateau failures persist.

\section{Implementation checklist from the derivation}

\begin{enumerate}
\item Construct $C_t=I_t^{\logi}$ by default, including dormant plateau
coordinates by provenance status.
\item Compute raw horizontal tangent Gram primitives $G_C,q_r,F_r$.
\item Compute $\sigma^\perp$, $\mathcal N^{\rm frac}$, and
$\Delta V^\lambda$.
\item Apply absolute and fractional residual gates before ranking.
\item Rank by $[\Delta V^\lambda]_{\lcb}/(1+K_3)$.
\item Log all geometry components, not only the final score.
\item Run the ordinary enlarged refit over active, dormant, and selected
coordinates.
\item Commit nonzero amplitudes only after validated unlock; otherwise commit
at zero and restore optimizer state.
\end{enumerate}

\section{One-sentence doctrine}

During a plateau, Route C should use Fubini--Study metric geometry to buy the
candidate-position record that most increases the resolved local log-volume of
the admitted active-plus-dormant scaffold, then let the ordinary enlarged refit
try to turn that purchased coordinate chart into a validated energy unlock.

\end{document}
